\section{方阵的行列式}

设 n 是一个正整数。我们把上述结果应用于体 D 上的右向量空间 \(D_{d}^{n}\)；\(D_{d}^{n}\) 的元素被视为 n 行一列的矩阵。设 \((\varepsilon_{1},\ldots,\varepsilon_{n})\) 是 \(D_{d}^{n}\) 的典范基。由 VIII, 第 451 页的推论 2，存在唯一的元素 \(\omega_{0}\) 属于 \(\Omega(\mathrm{D}_{d}^{n})\)，使得 \(\omega_{0}(\varepsilon_{1},\ldots,\varepsilon_{n})=1\)。若 A 是 \(\mathbf{GL}_{n}(\mathrm{D})\) 的一个元素，则它的列 \(a_{1},\ldots,a_{n}\) 构成 \(D_{d}^{n}\) 的一个基；元素 \(\omega_{0}(a_{1},\ldots,a_{n})\)（\(D_{ab}^{*}\) 中的）称为 A 的行列式，记作 \(\det(A)\)。由于我们有 \(a_{i}=A\varepsilon_{i}\)（对 \(1\leqslant i\leqslant n\)），A 的行列式就是自同构 \(x\mapsto Ax\)（\(D_{d}^{n}\) 的）的行列式。特别地，若体 D 交换，则 A 的行列式与 III, §8, No.~3, 第 524 页中定义的一致。

设 \(\mathrm{V}\) 是一个维数 \(n\) 有限的右向量空间（体 \(\mathrm{D}\) 上的），且 \((e_1, \ldots, e_n)\) 是 \(\mathrm{V}\) 的一个基。若 \(u\) 是 \(\mathrm{V}\) 的一个自同构，且 \(A\) 是 \(u\) 关于基 \((e_1, \ldots, e_n)\) 的矩阵，则我们有 \(\det(u) = \det(A)\)。

用 \((E_{ij})\) 表示 \(\mathbf{M}_{n}(\mathrm{D})\) 的矩阵单位族（II, §10, No.~3, 第 341 页）。对 D 的每个元素 \(\lambda\) 与每对不同的整数 \((i,j)\)（在 {[}1, n{]} 中），我们置（II, §10, No.~13, 第 361 页）

\[
B _ {i j} (\lambda) = I _ {n} + \lambda E _ {i j};
\]

它是 \(\mathbf{GL}_{n}(\mathrm{D})\) 的一个元素。若 \(\lambda_{1},\ldots,\lambda_{n}\) 是 \(D^{*}\) 的元素，则对角矩阵 \(\mathrm{diag}(\lambda_{1},\ldots,\lambda_{n})\) 也属于 \(\mathbf{GL}_{n}(\mathrm{D})\)。

命题 3. --- 映射 det 是从 \(\mathbf{GL}_{n}(\mathrm{D})\) 到 \(D_{ab}^{*}\) 的满足关系式

\[
\det (B _ {i j} (1)) = 1\tag{16}
\]

（对 \(i\neq j\)）与

\[
\det (\mathrm{diag} (\lambda_ {1}, \ldots , \lambda_ {n})) = \pi (\lambda_ {1} \dots \lambda_ {n}),\tag{17}
\]

（对 \(\lambda_1, \ldots, \lambda_n \in \mathrm{D}^*\)）的唯一同态。

设 \(A\) 是 \(\mathbf{GL}_n(\mathrm{D})\) 的一个元素，\(a_1, \ldots, a_n\) 是它的列。我们有

\[
\det (A) = \omega_ {0} (a _ {1}, \dots , a _ {n}).
\]

现在，矩阵 \(A \operatorname{diag}(\lambda_1, \ldots, \lambda_n)\) 的列是 \(a_1\lambda_1, \ldots, a_n\lambda_n\)，而矩阵 \(AB_{ij}(1)\) 的列因而是 \(a_1, \ldots, a_j + a_i, a_{j+1}, \ldots, a_n\)。由于 \(\omega_0\) 是 \(\Omega(\mathrm{D}_d^n)\) 中使 \(\omega_0(\varepsilon_1, \ldots, \varepsilon_n) = 1\) 的唯一元素，我们看出，行列式是满足下列关系式的唯一映射 \(\varphi: \mathbf{GL}_n(\mathrm{D}) \to \mathrm{D}_{\mathrm{ab}}^*\)

\begin{enumerate}
\def\labelenumi{(\arabic{enumi})}
\setcounter{enumi}{17}
\tightlist
\item
\end{enumerate}

\[
\varphi (A \operatorname{diag} (\lambda_ {1}, \ldots , \lambda_ {n})) = \varphi (A) \pi (\lambda_ {1} \dots \lambda_ {n}),\tag{19}
\]

\[
\varphi (A B _ {i j} (1)) = \varphi (A) \quad \mathrm{for} i \neq j,\tag{20}
\]

\[
\varphi (I _ {n}) = 1.
\]

这首先表明行列式映射满足关系式 (16) 与 (17)。我们已经知道这个映射是从 \(\mathbf{GL}_{n}(\mathrm{D})\) 到 \(D_{ab}^{*}\) 的一个同态。反之，若 \(\varphi\) 是从 \(\mathbf{GL}_{n}(\mathrm{D})\) 到 \(D_{ab}^{*}\) 的一个同态，使得 \(\varphi(B_{ij}(1)) = 1\)（对 \(i \neq j\)），且 \(\varphi(\text{diag}(\lambda_1, \ldots, \lambda_n)) = \pi(\lambda_1 \cdots \lambda_n)\)，则 \(\varphi\) 满足关系式 (18) 至 (20)，从而等于 det。

例. --- 1) 矩阵 \(B_{ij}(\lambda)\) 的列是 \(\varepsilon_1, \ldots, \varepsilon_j + \varepsilon_i \lambda, \varepsilon_{i+1}, \ldots, \varepsilon_n\)。由 VIII, 第 448 页的公式 (5)，我们有

\[
\det (B _ {i j} (\lambda)) = 1.\tag{21}
\]

\begin{enumerate}
\def\labelenumi{\arabic{enumi})}
\setcounter{enumi}{1}
\tightlist
\item
  设 \(\sigma\) 是区间 \([1,n]\)（在 \(\mathbf{N}\) 中）的一个置换，其符号差为 \(\varepsilon(\sigma)\)。设 \(M(\sigma)\) 是置换 \(\sigma\) 的矩阵（II, §10, No.~7, 第 351 页）。矩阵 \(M(\sigma)\) 的列是 \(\varepsilon_{\sigma(1)},\ldots,\varepsilon_{\sigma(n)}\)。于是应用 VIII, 第 448 页的公式 (4)，我们有
\end{enumerate}

\[
\det (M (\sigma)) = \pi (\varepsilon (\sigma)).\tag{22}
\]

\begin{enumerate}
\def\labelenumi{\arabic{enumi})}
\setcounter{enumi}{2}
\item
  设 \(n \geqslant 1\)。对每个形如 \(\Delta = \text{diag}(d_1, \ldots, d_n)\) 的可逆对角矩阵，我们有 \(\Delta B_{ij}(\lambda)\Delta^{-1} = B_{ij}(d_i\lambda d_j^{-1})\)。设 A 是 \(\mathbf{GL}_n(\mathrm{D})\) 的一个元素。由 II, §10, No.~13, 第 362 页的推论 1 以及前一个公式，存在矩阵 P 与 \(\Delta\)（在 \(\mathbf{GL}_n(\mathrm{D})\) 中），使得 \(A = P\Delta\)，P 是形如 \(B_{ij}(\lambda)\) 的矩阵的乘积，且 \(\Delta\) 是形如 \(\text{diag}(1, \ldots, 1, d)\) 的对角矩阵。由例 1 我们有 \(\det(P) = 1\)，从而由命题 3 有 \(\det(A) = \det(\Delta) = \pi(d)\)。
\item
  设 \(D'\) 是一个体，u 是从 D 到 \(D'\) 的一个同态。过渡到商群，u 定义一个群同态 \(u_{ab}\)（从 \(D_{ab}^{*}\) 到 \(D_{ab}^{\prime *}\)）。设 \(u_{n}\) 是从 \(\mathbf{GL}_{n}(\mathrm{D})\) 到 \(\mathbf{GL}_{n}(\mathrm{D}')\) 的把矩阵 \(A = (a_{ij})\) 变为矩阵 \((u(a_{ij}))\) 的同态。公式
\end{enumerate}

\[
\det (u _ {n} (A)) = u _ {\mathrm{ab}} (\det (A))\tag{23}
\]

对 \(A \in \mathbf{GL}_n(\mathrm{D})\) 立即由例 3 得出。

\begin{enumerate}
\def\labelenumi{\arabic{enumi})}
\setcounter{enumi}{4}
\item
  我们有 \(\mathbf{GL}_1(\mathrm{D}) = \mathrm{D}^*\)，且命题 3 表明我们有 \(\det(a) = \pi(a)\)（\(a\) 在 \(\mathbf{GL}_1(\mathrm{D})\) 中）。
\item
  设 \(n = 2\)。设 \(A = \begin{pmatrix} a & b \\ c & d \end{pmatrix}\) 是 \(\mathbf{GL}_2(\mathrm{D})\) 的一个元素。元素 \(a\) 与 \(c\)（\(\mathrm{D}\) 中的）不全为零。我们将明显地给出 \(A\) 的行列式。
\end{enumerate}

若 \(a\) 不为零，则我们有

\[
A = \left( \begin{array}{c c} 1 & 0 \\ c a ^ {- 1} & 1 \end{array} \right) \left( \begin{array}{c c} a & 0 \\ 0 & d - c a ^ {- 1} b \end{array} \right) \left( \begin{array}{c c} 1 & a ^ {- 1} b \\ 0 & 1 \end{array} \right).\tag{24}
\]

因此，\(ad - aca^{-1}b \neq 0\)，且

\[
\det (A) = \pi (a d - a c a ^ {- 1} b).\tag{25}
\]

\begin{enumerate}
\def\labelenumi{\alph{enumi})}
\setcounter{enumi}{1}
\tightlist
\item
  若 \(a\) 为零，则我们有 \(c \neq 0\)，且
\end{enumerate}

\[
A = \left( \begin{array}{c c} 0 & b \\ c & d \end{array} \right) = \left( \begin{array}{c c} 0 & 1 \\ 1 & 0 \end{array} \right) \left( \begin{array}{c c} c & d \\ 0 & b \end{array} \right).\tag{26}
\]

因此，由 a) 与例 2，我们有 \(cb \neq 0\)，且

\[
\det (A) = \pi (- c b).\tag{27}
\]

\begin{enumerate}
\def\labelenumi{\arabic{enumi})}
\setcounter{enumi}{6}
\tightlist
\item
  设 \(D^{o}\) 是 D 的反体。映射 \(A \mapsto {}^{t}A\)（从 \(\mathbf{M}_{n}(\mathrm{D})\) 到 \(\mathbf{M}_{n}(\mathrm{D}^{\circ})\)）满足 \({}^{t}(AB) = {}^{t}B^{t}A\)。因此，对 \(\mathbf{GL}_{n}(\mathrm{D})\) 中的每个矩阵 A，转置矩阵 \({}^{t}A\) 在 \(\mathbf{M}_{n}(\mathrm{D}^{\circ})\) 中可逆，但它在 \(\mathbf{M}_{n}(\mathrm{D})\) 中不一定可逆。由例 3 得出，若 A 属于 \(\mathbf{GL}_{n}(\mathrm{D})\)，则 \(\mathbf{GL}_{n}(\mathrm{D})\) 的元素 A 与 \({}^{t}A\)（\(\mathbf{GL}_{n}(\mathrm{D}^{\circ})\) 的）有相同的行列式。另一方面，即使矩阵 \({}^{t}A\) 属于 \(\mathbf{GL}_{n}(\mathrm{D})\)，它在 \(\mathrm{GL}_{n}(\mathrm{D})\) 中的行列式也不一定等于 A 的行列式（参见例 6）。
\end{enumerate}

注. --- 1) 设 V 是体 D 上的一个右向量空间，维数 n 有限。我们把对偶空间 \(V^{*} = \operatorname{Hom}_{\mathrm{D}}(\mathrm{V}, \mathrm{D})\) 视为 D 的反体 \(D^{o}\) 上的一个右向量空间。设 u 是 V 的一个自同构，\(^{t}u\) 是 \(V^{*}\) 的作为 u 的转置的自同构。若 u 由一个矩阵 A（在 \(\mathbf{M}_{n}(\mathrm{D})\) 中，关于 V 的一个基 \((e_{1}, \ldots, e_{n})\)）表示，则自同构 \(^{t}u\) 由矩阵 \(^{t}A\)（在 \(\mathbf{M}_{n}(\mathrm{D}^{o})\) 中）表示，所关于的基 \((e_{1}^{*}, \ldots, e_{n}^{*})\) 是 \(V^{*}\) 的与 \((e_{1}, \ldots, e_{n})\) 对偶的基（II, §10, No.~4, 第 344 页，命题 3）。由例 7，\(^{t}u\) 的行列式等于 u 的行列式。

\begin{enumerate}
\def\labelenumi{\arabic{enumi})}
\setcounter{enumi}{1}
\tightlist
\item
  第 1 至第 4 小节的结果可以推广到 D 是一个局部环的情形（VIII, 第 459 页，习题 2）。
\end{enumerate}

